% GENERATED FILE. Edit canonical lessons, then run npm run generate:books.
% canonical-source-hash: fnv1a64:3f78add2830c9e65
% canonical-lessons: ML-C475-kappal, ML-C475-navikan, ML-C475-jetti, ML-C475-vallam, ML-C475-tula

\chapter{Ships and Boats}
\label{ch:ml-ships-and-boats}
\hlchaptermodality{475}

\begin{chapteropening}
\textbf{By the end of this chapter:} \emph{I can say a ship, a sailor, a jetty, a country boat and an oar.}

Complete the last lesson of chapter 475: I can say a ship, a sailor, a jetty, a country boat and an oar.
\end{chapteropening}

\section[\texorpdfstring{\ml{കപ്പൽ} (kappal)}{കപ്പൽ (kappal)}]{\ml{കപ്പൽ} (kappal) — a ship}

\label{lesson:ML-C475-kappal}



\begin{quote}
Before the new one: say the Malayalam for to hang on, then the Malayalam for to speed along.
\end{quote}

\subsection*{You'll want to know: \ml{കപ്പൽ}}
\textbf{\ml{കപ്പൽ}} — \emph{kappal} — \textquotedblleft{}a ship\textquotedblright{}.

\begin{quote}
\ml{തുറമുഖത്ത്} \ml{ഒരു} \ml{വലിയ} \ml{കപ്പൽ} \ml{വന്നു}.
\end{quote}

\begin{quote}
\emph{tuṟamukhattŭ oru valiya kappal vannu.}
\end{quote}

\begin{quote}
A big ship came into the harbour.
\end{quote}

Plural: \ml{കപ്പലുകൾ}, \emph{kappalukaḷ}.

A small boat is a \ml{തോണി}; a \ml{കപ്പൽ} crosses the sea.

\begin{cousinweb}
\textbf{\ml{കപ്പൽ}} (\emph{kappal}) is native Dravidian, \textquotedblleft{}a ship, a sailing vessel\textquotedblright{} (DEDR 1219).
\end{cousinweb}

\begin{grammarlens}[title={What you've built}]
One more word for getting around.
\end{grammarlens}

\subsection*{Your turn}
\begin{itemize}
\raggedright
  \item \emph{Say it:} \emph{kappal}
  \item \emph{Say it:} \emph{kappal}, once more
  \item \emph{Say it:} \emph{tuṟamukhattŭ oru valiya kappal vannu.} — A big ship came into the harbour.
  \item \emph{Recall:} say the Malayalam for to hang on, then the Malayalam for to speed along, then say \emph{kappal} again
\end{itemize}

\begin{tcolorbox}[breakable,colback=teal!4,colframe=teal!35!black,title={Before you move on}]
What does \ml{കപ്പൽ} mean? (\textquotedblleft{}A ship\textquotedblright{}.) Say it once more.
\end{tcolorbox}
\section[\texorpdfstring{\ml{നാവികൻ} (nāvikan)}{നാവികൻ (nāvikan)}]{\ml{നാവികൻ} (nāvikan) — a sailor}

\label{lesson:ML-C475-navikan}



\begin{quote}
Before the new one: say the Malayalam for to speed along, then the Malayalam for a ship.
\end{quote}

\subsection*{You'll want to know: \ml{നാവികൻ}}
\textbf{\ml{നാവികൻ}} — \emph{nāvikan} — \textquotedblleft{}a sailor\textquotedblright{}.

\begin{quote}
\ml{നാവികൻ} \ml{കപ്പലിൽ} \ml{നിന്ന്} \ml{കരയിലേക്ക്} \ml{ഇറങ്ങി}.
\end{quote}

\begin{quote}
\emph{nāvikan kappalil ninnŭ karayilēkkŭ iṟaṅṅi.}
\end{quote}

\begin{quote}
The sailor came down from the ship onto the land.
\end{quote}

The navy is the \ml{നാവികസേന}, the sailors' army.

\begin{cousinweb}
\textbf{\ml{നാവികൻ}} (\emph{nāvikan}) is borrowed from Sanskrit \emph{nāvika}, \textquotedblleft{}belonging to a boat\textquotedblright{}, and so a boatman or sailor (Monier-Williams), with the masculine ending -\ml{ൻ}.
\end{cousinweb}

\begin{grammarlens}[title={What you've built}]
One more word for getting around.
\end{grammarlens}

\subsection*{Your turn}
\begin{itemize}
\raggedright
  \item \emph{Say it:} \emph{nāvikan}
  \item \emph{Say it:} \emph{nāvikan}, once more
  \item \emph{Say it:} \emph{nāvikan kappalil ninnŭ karayilēkkŭ iṟaṅṅi.} — The sailor came down from the ship onto the land.
  \item \emph{Recall:} say the Malayalam for to speed along, then the Malayalam for a ship, then say \emph{nāvikan} again
\end{itemize}

\begin{tcolorbox}[breakable,colback=teal!4,colframe=teal!35!black,title={Before you move on}]
What does \ml{നാവികൻ} mean? (\textquotedblleft{}A sailor\textquotedblright{}.) Say it once more.
\end{tcolorbox}
\section[\texorpdfstring{\ml{ജെട്ടി} (jeṭṭi)}{ജെട്ടി (jeṭṭi)}]{\ml{ജെട്ടി} (jeṭṭi) — a jetty}

\label{lesson:ML-C475-jetti}



\begin{quote}
Before the new one: say the Malayalam for a ship, then the Malayalam for a sailor.
\end{quote}

\subsection*{You'll want to know: \ml{ജെട്ടി}}
\textbf{\ml{ജെട്ടി}} — \emph{jeṭṭi} — \textquotedblleft{}a jetty\textquotedblright{}.

\begin{quote}
\ml{ഞങ്ങൾ} \ml{ജെട്ടിയിൽ} \ml{തോണിക്കായി} \ml{കാത്തുനിന്നു}.
\end{quote}

\begin{quote}
\emph{ñaṅṅaḷ jeṭṭiyil tōṇikkāyi kāttuninnu.}
\end{quote}

\begin{quote}
We stood waiting for the boat at the jetty.
\end{quote}

Boats on the backwaters stop at a \ml{ജെട്ടി} the way a bus stops at a stop.

\begin{cousinweb}
\textbf{\ml{ജെട്ടി}} (\emph{jeṭṭi}) was borrowed from English \emph{jetty}, a landing stage for boats.
\end{cousinweb}

\begin{grammarlens}[title={What you've built}]
One more word for getting around.
\end{grammarlens}

\subsection*{Your turn}
\begin{itemize}
\raggedright
  \item \emph{Say it:} \emph{jeṭṭi}
  \item \emph{Say it:} \emph{jeṭṭi}, once more
  \item \emph{Say it:} \emph{ñaṅṅaḷ jeṭṭiyil tōṇikkāyi kāttuninnu.} — We stood waiting for the boat at the jetty.
  \item \emph{Recall:} say the Malayalam for a ship, then the Malayalam for a sailor, then say \emph{jeṭṭi} again
\end{itemize}

\begin{tcolorbox}[breakable,colback=teal!4,colframe=teal!35!black,title={Before you move on}]
What does \ml{ജെട്ടി} mean? (\textquotedblleft{}A jetty\textquotedblright{}.) Say it once more.
\end{tcolorbox}
\section[\texorpdfstring{\ml{വള്ളം} (vaḷḷaṁ)}{വള്ളം (vaḷḷaṁ)}]{\ml{വള്ളം} (vaḷḷaṁ) — a country boat}

\label{lesson:ML-C475-vallam}



\begin{quote}
Before the new one: say the Malayalam for a sailor, then the Malayalam for a jetty.
\end{quote}

\subsection*{You'll want to know: \ml{വള്ളം}}
\textbf{\ml{വള്ളം}} — \emph{vaḷḷaṁ} — \textquotedblleft{}a country boat\textquotedblright{}.

\begin{quote}
\ml{രാവിലെ} \ml{ആ} \ml{വള്ളം} \ml{കായലിലൂടെ} \ml{പോയി}.
\end{quote}

\begin{quote}
\emph{rāvile ā vaḷḷaṁ kāyalilūṭe pōyi.}
\end{quote}

\begin{quote}
In the morning that boat went along the backwaters.
\end{quote}

You met it inside \ml{വള്ളംകളി}, the boat race. A boatman is a \ml{വള്ളക്കാരൻ}.

\begin{cousinweb}
\textbf{\ml{വള്ളം}} (\emph{vaḷḷaṁ}) is native Dravidian; DEDR gives the Malayalam word as a canoe hollowed from one trunk (DEDR 5315).
\end{cousinweb}

\begin{grammarlens}[title={What you've built}]
One more word for getting around.
\end{grammarlens}

\subsection*{Your turn}
\begin{itemize}
\raggedright
  \item \emph{Say it:} \emph{vaḷḷaṁ}
  \item \emph{Say it:} \emph{vaḷḷaṁ}, once more
  \item \emph{Say it:} \emph{rāvile ā vaḷḷaṁ kāyalilūṭe pōyi.} — In the morning that boat went along the backwaters.
  \item \emph{Recall:} say the Malayalam for a sailor, then the Malayalam for a jetty, then say \emph{vaḷḷaṁ} again
\end{itemize}

\begin{tcolorbox}[breakable,colback=teal!4,colframe=teal!35!black,title={Before you move on}]
What does \ml{വള്ളം} mean? (\textquotedblleft{}A country boat\textquotedblright{}.) Say it once more.
\end{tcolorbox}
\section[\texorpdfstring{\ml{തുഴ} (tuḻa)}{തുഴ (tuḻa)}]{\ml{തുഴ} (tuḻa) — an oar; a paddle}

\label{lesson:ML-C475-tula}



\begin{quote}
Before the new one: say the Malayalam for a jetty, then the Malayalam for a country boat.
\end{quote}

\subsection*{You'll want to know: \ml{തുഴ}}
\textbf{\ml{തുഴ}} — \emph{tuḻa} — \textquotedblleft{}an oar; a paddle\textquotedblright{}.

\begin{quote}
\ml{വള്ളക്കാരൻ} \ml{രണ്ട്} \ml{തുഴകൾ} \ml{എടുത്തു}.
\end{quote}

\begin{quote}
\emph{vaḷḷakkāran raṇṭŭ tuḻakaḷ eṭuttu.}
\end{quote}

\begin{quote}
The boatman picked up two oars.
\end{quote}

Plural: \ml{തുഴകൾ}, \emph{tuḻakaḷ}.

The verb is \ml{തുഴയുക}, to row: a \ml{തുഴ} is what you row with.

\begin{cousinweb}
\textbf{\ml{തുഴ}} (\emph{tuḻa}) is native Dravidian, the noun of the root of \ml{തുഴയുക}, to row or paddle (DEDR 3356).
\end{cousinweb}

\begin{grammarlens}[title={What you've built}]
One more word for getting around.
\end{grammarlens}

\subsection*{Your turn}
\begin{itemize}
\raggedright
  \item \emph{Say it:} \emph{tuḻa}
  \item \emph{Say it:} \emph{tuḻa}, once more
  \item \emph{Say it:} \emph{vaḷḷakkāran raṇṭŭ tuḻakaḷ eṭuttu.} — The boatman picked up two oars.
  \item \emph{Recall:} say the Malayalam for a jetty, then the Malayalam for a country boat, then say \emph{tuḻa} again
\end{itemize}

\begin{tcolorbox}[breakable,colback=teal!4,colframe=teal!35!black,title={Before you move on}]
What does \ml{തുഴ} mean? (\textquotedblleft{}An oar; a paddle\textquotedblright{}.) Say it once more.
\end{tcolorbox}
